\documentclass[11pt,a4paper]{article}
\usepackage[utf8]{inputenc}
\usepackage[T1]{fontenc}
\usepackage{amsmath,amsfonts,amssymb}
\usepackage{graphicx}
\usepackage{hyperref}
\usepackage{listings}
\usepackage{color}
\usepackage{booktabs}
\usepackage{float}
\usepackage{geometry}
\geometry{margin=1in}
\definecolor{zenblue}{RGB}{41,121,255}
\definecolor{zengreen}{RGB}{52,199,89}
\hypersetup{colorlinks=true,linkcolor=zenblue,urlcolor=zenblue,citecolor=zenblue}

\title{\textbf{Zen-VL: A Packaging of the Qwen3-VL Series}\\
\large Technical Note v2025.06}
\author{Zen LM Research Team\\
\texttt{research@zenlm.org}}
\date{June 2025}

\begin{document}
\maketitle

\begin{abstract}
Zen-VL is \emph{not} a from-scratch model. It is a redistribution and packaging of the
\textbf{Qwen3-VL} vision-language model series developed by the Qwen team at Alibaba Cloud
and released under the Apache-2.0 license~\cite{qwen3vl}. The Zen-VL packaging targets the
compact dense editions of that series --- \textbf{Qwen3-VL-4B} and \textbf{Qwen3-VL-8B}
(each available in Instruct and reasoning-oriented ``Thinking'' editions) --- with the
larger Mixture-of-Experts variants (30B-A3B and 235B-A22B) available upstream for
higher-capacity deployments. This note documents the upstream architecture as published by
its authors, records the provenance and license obligations, and describes the thin
packaging layer (weight conversion, quantization, serving configuration) Zen LM applies. No
architectural novelty, no separate pre-training run, and no independently produced benchmark
results are claimed. For authoritative capability numbers, see the upstream model cards and
the Qwen3-VL Technical Report~\cite{qwen3vl,qwen3vl4b,qwen3vlreport}.
\end{abstract}

\tableofcontents
\newpage

%% ─────────────────────────────────────────────────────────────────────────────
\section{Provenance and Scope}

Earlier internal drafts of this document described a homegrown 7B ``Zen MoDE (Mixture of
Distilled Experts)'' backbone with a ViT-L/14 encoder, a bespoke dynamic-tiling scheme,
a two-stage training recipe, and a set of benchmark tables presented as in-house results.
That framing was inaccurate and has been removed. The artifact distributed as ``Zen-VL'' is
a packaging of the Qwen3-VL series published by Alibaba Cloud's Qwen team~\cite{qwen3vl}.
This note exists to attribute the work correctly and to make the redistribution terms
explicit.

\textbf{Upstream models.} Qwen3-VL is the current-generation multimodal LLM series from the
Qwen team~\cite{qwen3vlreport}. It spans dense models at 4B and 8B parameters and MoE
variants at 30B-A3B and 235B-A22B, each offered in Instruct and ``Thinking''
editions~\cite{qwen3vl4b}. Zen-VL packages the dense \textbf{4B} and \textbf{8B} editions as
the compact entry points; the upstream MoE editions are documented separately and may be
packaged under the Zen3-VL note.

\textbf{License.} The upstream weights are released under \textbf{Apache-2.0}. This is
permissive and allows commercial use and redistribution, but a redistributor must preserve
copyright and license notices, include the upstream \texttt{NOTICE} file if present, and
state any modifications. Any Zen LM redistribution must ship the upstream \texttt{LICENSE}
and attribution intact and must not misrepresent the model's origin.

\textbf{Correction of model identity.} The ``7B Zen MoDE'' backbone, the ``307M ViT-L/14''
encoder, the specific tiling configuration table, the ten-million-pair Stage-1 and
two-million-pair Stage-2 training recipe, and all ``Zen-VL vs.\ Comp.\ A/B/C'' benchmark
tables in earlier drafts were fabricated and do not describe these models. The actual
artifacts are the Qwen3-VL-4B and Qwen3-VL-8B dense checkpoints.

\textbf{What ``packaging'' means here.} Zen LM does not retrain or re-architect the model.
The packaging layer is limited to weight-format conversion, optional post-training
quantization, and serving configuration.

%% ─────────────────────────────────────────────────────────────────────────────
\section{Upstream Architecture (as published by Qwen)}

The description below summarizes the architecture as documented by the upstream authors in
the Qwen3-VL model cards and Technical Report~\cite{qwen3vl4b,qwen3vlreport}. It is
reproduced for convenience and is \emph{not} a Zen LM contribution.

\subsection{Model Family}

\begin{table}[H]
\centering
\caption{Qwen3-VL series as published upstream (Alibaba Cloud, Apache-2.0). Zen-VL packages
the dense 4B/8B editions.}
\label{tab:family}
\begin{tabular}{lll}
\toprule
\textbf{Upstream model} & \textbf{Type} & \textbf{Editions} \\
\midrule
Qwen3-VL-4B        & Dense        & Instruct, Thinking \\
Qwen3-VL-8B        & Dense        & Instruct, Thinking \\
Qwen3-VL-30B-A3B   & MoE          & Instruct, Thinking \\
Qwen3-VL-235B-A22B & MoE          & Instruct, Thinking \\
\bottomrule
\end{tabular}
\end{table}

\subsection{Three-Module Structure}

Following the Qwen2.5-VL design, each Qwen3-VL model adopts a three-module structure: a
ViT-based vision encoder, an MLP-based vision--language merger that projects visual features
into the language embedding space, and the language model backbone (dense for the 4B/8B
editions). Exact layer counts, hidden sizes, and patch configurations are defined by the
upstream configuration files and are not restated or invented here.

\subsection{Named Architectural Features}

The upstream authors highlight several features shared across the Qwen3-VL
series~\cite{qwen3vl4b}:

\begin{itemize}
  \item \textbf{Interleaved-MRoPE} --- multimodal rotary position embeddings with
        full-frequency allocation over the time, width, and height axes.
  \item \textbf{DeepStack} --- fusion of multi-level ViT features to capture fine-grained
        visual detail and improve image--text alignment.
  \item \textbf{Text--Timestamp Alignment} --- timestamp-grounded event localization for
        stronger video temporal modeling.
\end{itemize}

\subsection{Context and Resolution}

Per the upstream model cards, the Qwen3-VL series provides a native 256K-token context
window, expandable to 1M tokens, with high-resolution image and video
input~\cite{qwen3vl4b}. Multi-image conversations and video understanding are supported
natively. Exact resolution handling is defined upstream.

\subsection{Items Removed From Earlier Drafts}

The following claims appeared in prior drafts and are \textbf{withdrawn} because they were
not substantiated and were presented as in-house work: the 7B ``Zen MoDE'' backbone and its
7.5B-total component table; the ViT-L/14 (307M) encoder spec and a 1024-token-per-tile
patching claim; the specific dynamic-tiling configuration and visual-token-count tables; the
Stage-1/Stage-2 training recipe with stated optimizer, batch size, and data counts; the SFT
data-composition table; and \emph{all} evaluation tables (multimodal benchmarks, document
understanding, high-resolution ablation, multi-image, and hallucination). These numbers were
invented. Genuine capabilities belong to the upstream Qwen3-VL models.

%% ─────────────────────────────────────────────────────────────────────────────
\section{Packaging and Redistribution}

\subsection{Conversion and Quantization}

Zen LM converts the published Qwen3-VL-4B and Qwen3-VL-8B weights to the internal serving
format and may apply post-training quantization for deployment, which is especially relevant
for the edge-oriented 4B edition. These transforms are mechanical; they do not constitute a
new training run and do not produce a new model. Where quantization is applied, the scheme
and bit-width should be recorded alongside the artifact so downstream users understand any
quality trade-off relative to the upstream checkpoint.

\subsection{Attribution Requirements}

Because the upstream license is Apache-2.0, every redistributed Zen-VL artifact must:

\begin{enumerate}
  \item Retain the upstream copyright and Apache-2.0 \texttt{LICENSE} text.
  \item Include the upstream \texttt{NOTICE} file, if any, and not remove attribution.
  \item State plainly that the model is a packaging of Qwen3-VL (4B/8B) by the Qwen team at
        Alibaba Cloud, and describe any modifications (e.g.\ quantization).
  \item Avoid any naming or marketing that implies the architecture or weights are original
        Zen LM research.
\end{enumerate}

%% ─────────────────────────────────────────────────────────────────────────────
\section{Evaluation: Use Upstream Numbers}

This note deliberately reports \textbf{no} Zen-produced benchmark scores. The Qwen team
publishes Qwen3-VL evaluations on standard vision-language suites --- including general
multimodal reasoning (e.g.\ MMMU, MMBench), document and OCR understanding (e.g.\ DocVQA,
OCRBench, InfoVQA), chart and diagram tasks (e.g.\ ChartQA, AI2D), and object-hallucination
evaluation (e.g.\ POPE, HallusionBench) --- in the official model cards and Technical
Report~\cite{qwen3vl4b,qwen3vlreport}.

\textbf{Guidance for downstream reporting.} Cite upstream-reported numbers with explicit
attribution to Qwen, and link to the model card for the specific size (4B or 8B). If Zen LM
applies quantization, any re-measured scores should be reported as ``Qwen3-VL-4B/8B (Zen
packaging, \textit{N}-bit)'' with the evaluation harness named, framed as a measurement of
the upstream model under a deployment configuration --- never as the result of a distinct
model.

%% ─────────────────────────────────────────────────────────────────────────────
\section{Related Work}

The relevant prior art is the upstream lineage itself: the Qwen-VL and Qwen2.5-VL line of
vision-language models from Alibaba Cloud, of which Qwen3-VL is the current
generation~\cite{qwen3vl4b,qwen3vlreport}. Contrastive vision-language
pre-training~\cite{radford2021clip} and projector-based connection of frozen vision encoders
to language models (BLIP-2~\cite{li2023blip2}, LLaVA~\cite{liu2023llava}) are foundational to
the broader field; the Interleaved-MRoPE, DeepStack, and Text--Timestamp Alignment
mechanisms used here are Qwen3-VL contributions.

%% ─────────────────────────────────────────────────────────────────────────────
\section{Limitations}

The limitations of Zen-VL are the limitations of the upstream Qwen3-VL models it packages,
as documented by their authors. Additionally, any post-training quantization applied during
packaging may reduce quality relative to the upstream full-precision checkpoints; the degree
depends on the chosen scheme and bit-width and should be measured and disclosed per artifact.

%% ─────────────────────────────────────────────────────────────────────────────
\section{Conclusion}

``Zen-VL'' denotes a packaged redistribution of Alibaba Cloud's Apache-2.0 Qwen3-VL dense
models (4B and 8B), not an independent model. The honest description of the work is:
convert, optionally quantize, serve, and attribute. All architectural credit and all
capability claims belong to the upstream Qwen3-VL authors and should be cited to them.

\section*{Acknowledgments}

We acknowledge the Qwen team at Alibaba Cloud as the authors of the Qwen3-VL series, the
models packaged and redistributed here under the Apache-2.0 license.

%% ─────────────────────────────────────────────────────────────────────────────
\begin{thebibliography}{99}

\bibitem{qwen3vl}
Qwen Team, Alibaba Cloud,
``Qwen3-VL,'' GitHub repository, 2025.
\url{https://github.com/QwenLM/Qwen3-VL}

\bibitem{qwen3vl4b}
Qwen Team, Alibaba Cloud,
``Qwen3-VL-4B-Instruct,'' Hugging Face model card, 2025.
\url{https://huggingface.co/Qwen/Qwen3-VL-4B-Instruct}

\bibitem{qwen3vlreport}
Qwen Team, Alibaba Cloud,
``Qwen3-VL Technical Report,'' arXiv:2511.21631, 2025.
\url{https://arxiv.org/abs/2511.21631}

\bibitem{radford2021clip}
A.~Radford et al., ``Learning Transferable Visual Models From Natural Language Supervision,''
\textit{ICML}, 2021.

\bibitem{li2023blip2}
J.~Li et al., ``BLIP-2: Bootstrapping Language-Image Pre-training with Frozen Image Encoders
and Large Language Models,'' \textit{ICML}, 2023.

\bibitem{liu2023llava}
H.~Liu et al., ``Visual Instruction Tuning,'' \textit{NeurIPS}, 2023.

\bibitem{dosovitskiy2020vit}
A.~Dosovitskiy et al., ``An Image is Worth 16x16 Words: Transformers for Image Recognition
at Scale,'' \textit{ICLR}, 2021.

\end{thebibliography}

\end{document}
